\section*{Chapter \ref{ch:ll:cpp-for-latency-iii-the-compiler} --- C++ for Latency III: The Compiler}

\begin{solution}{exo:ll:cpp-for-latency-iii-the-compiler:1}
On the committed measurement about 92, 11.6 and \qty{3.1}{\nano\second} per message; at 5 million messages a second, 460,
58 and 16 milliseconds of processing a second: the unoptimised build would use half a core on decoding alone.
\end{solution}

\begin{solution}{exo:ll:cpp-for-latency-iii-the-compiler:2}
\texttt{INT\_MAX + 1} is undefined (signed overflow). \texttt{UINT\_MAX + 1u} is defined: unsigned arithmetic wraps to 0.
\texttt{std::memcpy} into a local is defined. The cast is undefined: it reads an object of a type the buffer never held
(aliasing), and the address may be misaligned.
\end{solution}

\begin{solution}{exo:ll:cpp-for-latency-iii-the-compiler:3}
The comparison is folded while the compiler simplifies expressions, before any optimisation level applies: with signed
overflow impossible, \texttt{x + 1 > x} is always true. No addition is left in the program, so the sanitiser has nothing
to instrument; it can only catch an overflow that is actually computed, as in \texttt{increment}.
\end{solution}

\begin{solution}{exo:ll:cpp-for-latency-iii-the-compiler:4}
The exact sum is 14.440160. The \texttt{-ffast-math} result, 14.439274, is off by about 0.0009; the in-order sum, 14.403684,
by 0.036, forty times more, because adding tiny terms to a large running sum loses them. Closer is not safe: the result now
depends on the compiler, its version and the vector width, so two builds disagree, and the program can no longer be checked
against a reference. Where accuracy matters, write the better algorithm (pairwise or compensated summation, One Quant
Book~4, chapter~25) explicitly.
\end{solution}

\begin{solution}{exo:ll:cpp-for-latency-iii-the-compiler:5}
It breaks the aliasing rule (reading bytes as a \texttt{uint32\_t} object that does not exist) and the alignment rule
(\texttt{buf + 16} need not be 4-byte aligned), and it ignores the wire's byte order. Correct: \texttt{std::uint32\_t v;
std::memcpy(\&v, buf + 16, 4); v = \_\_builtin\_bswap32(v);} for a big-endian field on x86; it compiles to a load and a
byte swap.
\end{solution}

\begin{solution}{exo:ll:cpp-for-latency-iii-the-compiler:6}
The profile tells the compiler that rebalance-day paths (bursts of executions, auction messages, one-sided flow) are cold,
so it lays them out of line, does not inline them and may even optimise them for size: the day the firm most needs speed,
the hot paths are the ones built as cold. Train on representative busy days, including the rare heavy ones.
\end{solution}

\begin{solution}{exo:ll:cpp-for-latency-iii-the-compiler:7}
In \texttt{alias} the store through \texttt{f} is followed by a reload of \texttt{*i}
(\texttt{mov eax, DWORD PTR [rdi]}) instead of the constant 1. The decoder assembles fields byte by byte and does not rely
on type-based aliasing, so its code and speed should not change; the matrix measures it.
\end{solution}

\begin{solution}{exo:ll:cpp-for-latency-iii-the-compiler:8}
Three flaws: \texttt{-Ofast} includes \texttt{-ffast-math}, which changes floating-point results (the float sum changed); a
single benchmark number with no checksum proves speed, not correctness; and \texttt{-march=native} on the build server
targets the build server's processor, which may have instructions the trading hosts lack. Name the target, keep IEEE
semantics where results matter, and run the flag matrix with checksums on a production-shaped workload.
\end{solution}

\begin{solution}{pb:ll:cpp-for-latency-iii-the-compiler:1}
\begin{enumerate}
\item \texttt{-O0} 0.13; \texttt{-O3} 3.7; \texttt{-O2 -march=native} 1.02; \texttt{-O2 -flto} 1.02; \texttt{-O2} with
PGO 2.6; everything together 4.8 (committed measurement).
\item \texttt{-O3} alone, with \omterm{def:ll:cpp-for-latency-iii-the-compiler:pgo}{profile-guided optimisation} close behind.
\item The handler in the second file is cheap; the time is in the decoder and the per-message callback it makes, which are
in the main file already.
\item No: every set produced the \texttt{-O2} checksum.
\item \texttt{-O2}: \qty{0.44}{\nano\second} per element and 14.4036837; \texttt{-O3 -march=native}: the same time and
value; with \texttt{-ffast-math}: \qty{0.064}{\nano\second}, seven times faster, and 14.4392738.
\item It allowed the compiler to reassociate the additions, summing several partial sums in vector registers, a different
order with different rounding.
\item The \texttt{-ffast-math} value (exact 14.440160).
\item Its results must be reproducible and checkable against a reference, and \texttt{-ffast-math} also assumes no NaNs or
infinities, which market data and failed calibrations produce.
\item \omterm{def:ll:cpp-for-latency-iii-the-compiler:ub}{Undefined behaviour} gives no guarantee at any level: \texttt{-O3} inlines and vectorises more, so an assumption that
was harmless in one context becomes exploited in another.
\item ``runtime error: signed integer overflow: 2147483647 + 1 cannot be represented in type int''.
\item In test and continuous-integration builds, on replayed production days; not in production binaries.
\item That a check after a use is dead code to the compiler: check before use, and treat compiler warnings and sanitiser
reports as defects.
\item \textbf{Named result.} With \texttt{-O3}, \texttt{-march=native}, link-time and \omterm{def:ll:cpp-for-latency-iii-the-compiler:pgo}{profile-guided optimisation} the decoder
ran 4.8 times faster than at \texttt{-O2} with an identical checksum; \texttt{-ffast-math} made the float sum seven times
faster and changed its value.
\item \texttt{-O3} with \omterm{def:ll:cpp-for-latency-iii-the-compiler:pgo}{profile-guided optimisation}, an explicit target instead of \texttt{native}, no
\texttt{-ffast-math}; accompanied by a pinned compiler version, a training workload under version control, and the matrix
rerun on every release.
\item A replay of recent busy days, including the heaviest seen, distinct from the day used to measure.
\item To run only on processors with the build machine's instruction set; a new server generation may be fine, an older
one crashes with an illegal instruction.
\item It would compile the few hot vector functions for several targets and choose at load time, so the binary runs
everywhere and still uses AVX2 where available.
\item On every compiler upgrade, flag change or significant code change, and periodically on the current workload.
\item Because the bytes change with the compiler version while the properties (a call inlined, a check present) are what
the design relies on.
\item Nothing: faster and wrong is wrong.
\end{enumerate}
\end{solution}

\begin{solution}{iq:ll:cpp-for-latency-iii-the-compiler:1}
Behaviour on which the standard imposes no requirement (signed overflow, invalid pointers, data races, aliasing
violations). It exists so that compilers can assume it away and generate fast code for many machines without checks.
\iqlookfor{that the compiler optimises on the assumption it never happens.}
\end{solution}

\begin{solution}{iq:ll:cpp-for-latency-iii-the-compiler:2}
The compiler may assume that pointers to unrelated types never refer to the same memory. Read with \texttt{std::memcpy}
into a local (or \texttt{std::bit\_cast} on an array), then convert the byte order; it compiles to a load.
\iqlookfor{memcpy or bit\_cast, and byte order.}
\end{solution}

\begin{solution}{iq:ll:cpp-for-latency-iii-the-compiler:3}
It builds an instrumented binary, runs it on a training workload to record branch and value frequencies, and rebuilds
using them for \omterm{def:ll:cpp-for-latency-ii-compile-time:inline}{inlining}, layout and unrolling. Risks: a non-representative training run optimises the wrong paths; the
build depends on data that must be versioned.
\iqlookfor{the training workload as the critical input.}
\end{solution}

\begin{solution}{iq:ll:cpp-for-latency-iii-the-compiler:4}
Not for pricing, risk or anything checked against a reference; possibly for a self-contained numerical kernel whose
tolerance is tested, where the speed matters and NaNs cannot occur. Better: write the vectorised or reassociated algorithm
explicitly and keep IEEE semantics.
\iqlookfor{reproducibility and the NaN assumption.}
\end{solution}

\begin{solution}{iq:ll:cpp-for-latency-iii-the-compiler:5}
The pointer was dereferenced earlier; since dereferencing null is undefined, the compiler concludes it is not null and
removes the test. Move the check before the first use.
\iqlookfor{the use-then-check pattern.}
\end{solution}

\begin{solution}{iq:ll:cpp-for-latency-iii-the-compiler:6}
A flag matrix on a realistic workload with a checksum, pinned and repeated, including profile-guided builds trained on
separate data; reject any set that changes the output; read the hot assembly; pin the compiler and flags; rerun on every
upgrade and release.
\iqlookfor{measurement plus correctness, repeated over time.}
\end{solution}
